(1) 定态方程
\[
-\frac{\hbar^2}{2m}\psi''+mgz\psi=E\psi.
\]
定义长度 \(\ell_g=(\hbar^2/(2m^2g))^{1/3}\)、无量纲 \(\xi=z/\ell_g-E/(mg\ell_g)\)，方程成为 \(\psi_{\xi\xi}=\xi\psi\)，即 Airy 方程。

(2) 一般解由 \(\mathrm{Ai}(\xi),\mathrm{Bi}(\xi)\) 组成，但 \(\mathrm{Bi}(\xi)\) 在 \(\xi\to+\infty\) 发散，须剔除。墙面条件要求 \(\mathrm{Ai}(-\epsilon_n)=0\)，其中 \(\epsilon_n>0\) 是 Ai 的负零点绝对值。故
\[
E_n=mg\ell_g\epsilon_n,\quad
\psi_n(z)=N_n\mathrm{Ai}(z/\ell_g-\epsilon_n).
\]
由 Airy 正交积分 \(\int_0^\infty\mathrm{Ai}(x-\epsilon_n)^2dx=\mathrm{Ai}'(-\epsilon_n)^2\)，可取 \(N_n=[\ell_g^{1/2}|\mathrm{Ai}'(-\epsilon_n)|]^{-1}\)。

(3) 第一根 \(\epsilon_1\simeq2.3381074105\)，所以 \(E_1\simeq2.3381074105\,mg\ell_g\)，体现量子重力能量尺度 \((\hbar^2mg^2)^{1/3}\) 的幂律。

(4) 对线性势 \(V=mgz\)，量子维里关系 \(2\langle K\rangle=\langle z V'(z)\rangle=\langle V\rangle\)；硬墙虽提供约束，但因 \(\psi(0)=0\)、墙位于 z=0，维里边界项对该形式为零。因此 \(E=\langle K\rangle+\langle V\rangle=3\langle K\rangle\)，故
\[
\langle K\rangle=E_n/3,\qquad\langle V\rangle=2E_n/3.
\]
